% lecture12.tex: Lecture 12: Meteorites, Asteroids, Minor Planets, Comets
% Planetary Systems, Kapteyn Institute

\documentclass[aspectratio=169, 11pt]{beamer}

% ── Theme and macros ────────────────────────────────────────
\usepackage{../common/beamerthemeIPS}
% macros.tex: Shared math macros for IPS lecture slides
% Mirrors definitions in book/_config.yml for consistency
% Uses \providecommand to avoid conflicts with unicode-math

% ── Derivatives ─────────────────────────────────────────────
\providecommand{\dv}[2]{\frac{\mathrm{d} #1}{\mathrm{d} #2}}
\providecommand{\dd}{}\renewcommand{\dd}{\mathrm{d}}
\providecommand{\pdv}[2]{\frac{\partial #1}{\partial #2}}

% ── Solar quantities ────────────────────────────────────────
\newcommand{\Msun}{M_{\odot}}
\newcommand{\Rsun}{R_{\odot}}
\newcommand{\Lsun}{L_{\odot}}

% ── Planetary quantities ────────────────────────────────────
\newcommand{\Mearth}{M_{\oplus}}
\newcommand{\Rearth}{R_{\oplus}}
\newcommand{\Mjup}{M_{\mathrm{J}}}
\newcommand{\Rjup}{R_{\mathrm{J}}}

% ── Constants ───────────────────────────────────────────────
\newcommand{\kB}{k_{\mathrm{B}}}

% ── Media links ─────────────────────────────────────────────
% Clickable button that opens the animated or video version of a
% figure, e.g. \mediabutton{https://...gif}{Play animation}
\providecommand{\mediabutton}[2]{\href{#1}{\beamergotobutton{#2}}}

\graphicspath{{figures/}{../common/}{../../book/12_small_bodies/figures/}}

% ── Metadata ────────────────────────────────────────────────
\title[Lecture 12: Small Bodies]{Lecture 12: Meteorites, Asteroids,\\Minor Planets \& Comets}
\subtitle{Planetary Systems}
\author[L12]{Tim Lichtenberg}
\institute{Kapteyn Astronomical Institute\\University of Groningen}
\date{2026}
\titlebgcredit{Comet 1P/Halley, Giotto Halley Multicolour Camera (1986). Credit: ESA}

% ════════════════════════════════════════════════════════════
\begin{document}

% ── Title slide ─────────────────────────────────────────────
\begin{frame}[plain,noframenumbering]
  \titlepage
\end{frame}

% ── Learning objectives ─────────────────────────────────────
\begin{frame}{Learning Objectives}
  By the end of this lecture, you will be able to:
  \vskip8pt
  \begin{itemize}
    \item Classify meteorites and derive the Pb-Pb isochron age of calcium-aluminium-rich inclusions (CAIs)
    \item Describe the asteroid-belt and trans-Neptunian populations and their dynamics
    \item Explain the origin, anatomy and volatile composition of comets, and read small bodies as fossils of planet formation
  \end{itemize}
\end{frame}

% ── Recap ───────────────────────────────────────────────────
\begin{frame}{Recap: Where We Are in the Course}
  \begin{itemize}
    \item Lectures 9--11 surveyed the major planets across the rocky and giant classes
    \item Today: the leftovers; asteroids, meteorites, KBOs, comets, and the Oort cloud
    \item Small bodies act as formation fossils preserving disk chemistry and migration history
  \end{itemize}
  \vskip8pt
  \textbf{Today:} What do the leftovers reveal about the origin of planetary systems?
\end{frame}


% ════════════════════════════════════════════════════════════
\sectionimage{small_bodies_overview}{Size comparison of small bodies visited by spacecraft. NASA/JPL-Caltech}
\section{Overview: The Reservoirs}
% ════════════════════════════════════════════════════════════

\begin{frame}
  \ipsherokey[0.65]{small_bodies_overview}{Small bodies visited by spacecraft, scaled against terrestrial landmarks. Credit: NASA/JPL-Caltech, public domain.}{Three reservoirs hold them: the main belt ($2.1$--$3.3$~AU, $\sim 4 \times 10^{-4}\,\Mearth$), the Kuiper Belt, and the Oort cloud ($2{,}000$--$50{,}000$~AU, inferred, never directly observed).}
\end{frame}


% ════════════════════════════════════════════════════════════
\sectionimage{allende}{Allende meteorite (CV3 chondrite, fell 1969). J.\ St.\ John CC BY 2.0}
\section{Meteorite Taxonomy}
% ════════════════════════════════════════════════════════════

\begin{frame}{Why Meteorites Matter}
  \begin{itemize}\setlength\itemsep{4pt}
    \item \textbf{Laboratory access:} $\sim$70{,}000 catalogued specimens, older than the oldest terrestrial zircons ($\sim 4.4$~Ga)
    \item \textbf{Presolar material:} primitive chondrites contain grains that condensed around other stars, older than the Sun
    \item \textbf{Anchor age:} CAI crystallization age of $4567.30 \pm 0.16$~Myr anchors solar system history
  \end{itemize}
\end{frame}

\begin{frame}
  \ipsherokey[0.66]{allende}{Allende slab: chondrules and bright CAIs in a dark matrix. Credit: James St.\ John, CC BY 2.0.}{Chondrites never melted whole and keep their mm-scale chondrules; achondrites, irons and stony irons come from parent bodies that short-lived $^{26}\mathrm{Al}$ melted and differentiated: irons are cores, achondrites crusts and mantles.}
\end{frame}

\begin{frame}
  \ipsherokey[0.65]{chondrite_thin}{Cross-polarised thin section of NWA 5930: each chondrule is a mm-scale bead of olivine and pyroxene. Credit: H.\ Raab, CC BY-SA 3.0 (Wikimedia Commons).}{Each chondrule froze from a fully molten droplet in seconds to hours from a peak of $1500$--$1900$~K; how chondrules form remains an unsolved problem.}
\end{frame}

\begin{frame}
  \ipshero{chondrule_cai_timeline}{(a) CAIs at 4567.30 Myr define $t = 0$ and chondrules formed 2 to 4 Myr later; (b) a chondrule heating event from a 1700 K peak, cooling at 100 and 1000 K per hour. Course figure.}
\end{frame}

\begin{frame}
  \ipshero{petrologic_types}{Petrographic types 1 to 7 (aqueous alteration to thermal metamorphism) and shock stages S1 to S6 record what the parent body did. Course schematic.}
\end{frame}

\begin{frame}
  \ipsherokey[0.65]{iron_widmanstatten}{Widmanst\"atten pattern in an iron meteorite: a kamacite and taenite Fe-Ni lattice. Credit: H.\ Raab, CC BY-SA 3.0 (Wikimedia Commons).}{The lattice grows only at cooling rates $\lesssim 100$~K~Myr$^{-1}$, inside the metallic core of a body tens of km across; its spacing gives the original parent-body radius.}
\end{frame}

\begin{frame}
  \ipsherokey[0.65]{pallasite}{Esquel pallasite: yellow-green olivine in silvery Fe-Ni metal. Credit: UCLA Meteorite Gallery, CC BY-SA 4.0 (Wikimedia Commons).}{\textbf{A direct sample of a destroyed planetesimal's core-mantle boundary:} the olivine crystals sank into the molten metal during a giant impact.}
\end{frame}

\begin{frame}
  \ipsherokey[0.61]{vesta}{Vesta seen by Dawn (2011--2012), the parent body of the HED meteorites. Credit: NASA/JPL-Caltech/UCLA/MPS/DLR/IDA, public domain.}{HEDs (howardites, eucrites, diogenites, $\sim 6\%$ of falls) match Vesta in mineralogy and oxygen isotopes and were ejected from the Rheasilvia basin; lunar and Martian meteorites are identified the same way.}
\end{frame}

\begin{frame}
  \ipshero{oxygen_three_isotope}{Oxygen three-isotope plot: the terrestrial line of slope 1/2, the CAI slope-1 mixing line towards pure $^{16}$O, and the chondrite, HED and SNC fields. Course figure.}
\end{frame}


% ════════════════════════════════════════════════════════════
\sectionimage{pb_pb_dating}{Pb-Pb isochrons for CAI 22E and chondrules. Connelly et al.\ (2012)}
\section{Radiogenic Chronometers}
% ════════════════════════════════════════════════════════════

\begin{frame}{Decay Clocks and Early Solar System Timeline}
  Radioactive decay $\dv{N}{t} = -\lambda N$ ($t_{1/2} = \ln 2 / \lambda$) provides absolute and relative chronometers:
  \vskip4pt
  \begin{itemize}\setlength\itemsep{3pt}
    \item \textbf{Long-lived:} $^{238}\mathrm{U}$, $^{235}\mathrm{U}$ yield absolute ages; CAIs condense at $t_0 = 4567.30 \pm 0.16$~Myr
    \item \textbf{Short-lived (extinct):} $^{26}\mathrm{Al}$ ($0.717$~Myr), $^{182}\mathrm{Hf}$ ($8.9$~Myr) yield relative ages
    \item \textbf{Early chronology:} differentiation at $0.5$--$4$~Myr, chondrules at $1$--$4$~Myr after CAIs
  \end{itemize}
  \vskip6pt
  \keyresult{Long-lived chronometers anchor absolute time; extinct nuclides resolve events to $<1$~Myr.}
\end{frame}

\begin{frame}
  \ipsherokey{short_lived_decay}{Abundance decay of $^{26}$Al (half-life 0.72 Myr) and $^{60}$Fe (2.62 Myr) over the first 10 Myr after CAI formation. Course figure.}{Their rapid decay powered early planetesimal melting and differentiation.}
\end{frame}


% ════════════════════════════════════════════════════════════
\sectionimage{allende}{Allende slab. James St.\ John CC BY 2.0}
\section{Blackboard Derivation: Pb-Pb Isochron of CAIs}
% ════════════════════════════════════════════════════════════

\begin{frame}{Goal and Strategy}
  \begin{columns}[c]
    \column{\recapfigcolnarrow}
    \ipsrecapfig[0.62\textheight]{board_sketch_l12}
    \column{\recaptextcolwide}
    \small
    \textbf{Goal:} the age of the oldest solids from two parallel decay chains of one element.
    \vskip4pt
    \begin{itemize}\setlength\itemsep{2pt}
      \item $^{238}\mathrm{U} \to {}^{206}\mathrm{Pb}$ (4.4683~Gyr) and $^{235}\mathrm{U} \to {}^{207}\mathrm{Pb}$ (0.7038~Gyr); $^{204}\mathrm{Pb}$ is the stable reference
      \item Daughter since closure: ${}^{206}\mathrm{Pb}^* = {}^{238}\mathrm{U}\,[e^{\lambda_{238} t} - 1]$, same for 207; add the initial Pb
      \item Divide the 207 line by the 206 line so the U/Pb of the rock drops out
    \end{itemize}
    \vskip4pt
    \keyresult{To the board.}
  \end{columns}
\end{frame}

\begin{frame}{Result: The Pb-Pb Isochron}
  Cogenetic samples with unknown initial Pb fall on a line in $({}^{206}\mathrm{Pb}/{}^{204}\mathrm{Pb},\ {}^{207}\mathrm{Pb}/{}^{204}\mathrm{Pb})$ whose slope depends only on $t$:
  \[
    \boxed{\; \frac{{}^{207}\mathrm{Pb}^*}{{}^{206}\mathrm{Pb}^*} = \frac{{}^{235}\mathrm{U}}{{}^{238}\mathrm{U}}\,\frac{e^{\lambda_{235} t} - 1}{e^{\lambda_{238} t} - 1}, \qquad \frac{{}^{235}\mathrm{U}}{{}^{238}\mathrm{U}} = \frac{1}{137.8} \;}
  \]
  \keyresult{The slope is independent of the initial Pb and of the U/Pb ratio: the age needs no assumption about the starting composition.}
\end{frame}

\begin{frame}
  \ipshero{pb_pb_dating}{Pb-Pb isochrons for CAI 22E (Efremovka) and chondrules: cogenetic mineral fractions define a line whose slope gives the age $t$ independent of the initial Pb. From Connelly et al.\ (2012).}
\end{frame}

\begin{frame}{Step 5: Numerical Result}
  Connelly et al.\ (2012), CAIs from NWA 2364 (CV3) after acid leaching:
  \[
    \boxed{\; t_{\text{CAI}} = 4567.30 \pm 0.16\ \mathrm{Myr} \;}
  \]
  \begin{itemize}
    \item Independent samples and labs agree to the same precision (Amelin et al.\ 2010)
    \item Chondrules from the same meteorite: a few Myr younger
    \item The \emph{zero point} of the solar system clock
  \end{itemize}
  \vskip4pt
  \keyresult{Two parallel decay chains make the problem self-calibrating.}
\end{frame}


% ════════════════════════════════════════════════════════════
\sectionimage{spitzer_nc_cc_ages}{Hf-W ages of NC and CC iron meteorites. Spitzer et al.\ (2021)}
\section{The NC-CC Dichotomy}
% ════════════════════════════════════════════════════════════

\begin{frame}{Two Isotopic Reservoirs}
  Nucleosynthetic anomalies in Ti, Cr, Mo, Ru and Ca divide solar system materials into two isolated reservoirs:
  \vskip8pt
  \begin{itemize}\setlength\itemsep{4pt}
    \item \textbf{NC reservoir (non-carbonaceous):} ordinary and enstatite chondrites, HEDs, Earth, and Mars
    \item \textbf{CC reservoir (carbonaceous):} carbonaceous chondrites, several iron groups, and outer bodies
    \item \textbf{Prolonged isolation:} the two reservoirs did not mix isotopically for $\sim 3$--$4$~Myr after CAIs
  \end{itemize}
\end{frame}

\begin{frame}
  \ipsherokey[0.62]{spitzer_nc_cc_ages}{Hf-W ages of NC (red) and CC (blue) iron meteorites: $\varepsilon^{182}\mathrm{W}$ below, the age it gives above. Spitzer et al.\ (2021).}{NC irons ($-3.37$ to $-3.20$) formed cores $1.0$--$2.6$~Myr after CAIs, CC irons ($-3.16$ to $-3.10$) at $3.0$--$3.6$~Myr, the ungrouped SBT $\sim 1$~Myr earlier: NC cores form $\sim 2$~Myr before CC cores.}
\end{frame}

\begin{frame}{Three Live Interpretations}
  Three competing mechanisms explain the NC-CC dichotomy from identical cosmochemical data:
  \vskip8pt
  \begin{itemize}\setlength\itemsep{4pt}
    \item \textbf{Jupiter barrier:} early proto-Jupiter ($\lesssim 1$~Myr) opens a gap, blocking inward pebble drift (Kruijer et al.\ 2017)
    \item \textbf{Snow line:} disk thermal evolution produces two successive planetesimal formation epochs (Lichtenberg et al.\ 2021)
    \item \textbf{Temporal change:} inward disk drift alters feeding composition over time (Schiller et al.\ 2018)
  \end{itemize}
\end{frame}

\begin{frame}
  \ipsherokey[0.65]{lichtenberg2021_fig1}{Disk simulation in time against orbital distance with two planetesimal formation reservoirs. Lichtenberg et al.\ (2021).}{Reservoir I (red), early and dry in the inner disk, is NC; Reservoir II (blue), later and ice-rich in the outer disk, is CC: the bifurcation arises from disk evolution itself.}
\end{frame}


% ════════════════════════════════════════════════════════════
\sectionimage{demeo_carry_belt_inclination}{Asteroid orbits in $i$ vs $a$. DeMeo \& Carry (2014)}
\section{Asteroid Belt: Structure and Dynamics}
% ════════════════════════════════════════════════════════════

\begin{frame}
  \ipsherokey[0.65]{demeo_carry_belt_inclination}{Asteroid orbits in inclination against semimajor axis: Hungarias, the main belt at 2.1--3.3 AU, Hildas, and the Trojans at 5.2 AU. DeMeo \& Carry (2014).}{About 1.4 million catalogued bodies, $\sim 10^6$ above 1 km, a total mass of $\sim 4 \times 10^{-4}\,\Mearth$ (Ceres $\approx 40\%$); the vertical gaps are the Kirkwood resonances with Jupiter.}
\end{frame}

\begin{frame}
  \ipsherokey{granvik_kirkwood_resonances}{The Kirkwood gaps at the 3:1 (2.50 AU), 5:2 (2.82), 7:3 (2.96) and 2:1 (3.27) resonances with Jupiter. Granvik et al.\ (2018).}{Periodic Jovian kicks add coherently and pump eccentricities to planet-crossing in $10^4$--$10^6$~yr: the gaps are not empty but leaky, the dynamical sinks that feed the near-Earth population.}
\end{frame}

\begin{frame}
  \ipsherokey[0.65]{asteroid_families_nesvorny}{Asteroid families as clusters in proper $a$, $e$, $i$: Themis, Eos, Koronis, Eunomia, Vesta, Flora. Nesvorn\'y et al.\ (2015).}{Each family is the debris of a single disruption event; its spread gives the original parent body's size and age.}
\end{frame}

\begin{frame}
  \ipsherokey[0.61]{demeo_carry_belt_mass}{Mass of the main belt by spectral class against heliocentric distance. DeMeo \& Carry (2014).}{\textbf{S-types} ($\sim 15\%$ albedo, ordinary-chondrite analogues) fill the inner belt, \textbf{C-types} ($\sim 4\%$ albedo, volatile-rich) $\sim 75\%$ of the outer belt, \textbf{M- and V-types} are cores and Vesta's basaltic crust ($2.4$~AU): radial sorting preserves the disk's primordial volatile gradient.}
\end{frame}


% ════════════════════════════════════════════════════════════
\sectionimage{chelyabinsk}{Chelyabinsk meteor (2013). A.\ Alishevskikh CC BY-SA 2.0}
\section{NEAs, Yarkovsky/YORP, and Planetary Defence}
% ════════════════════════════════════════════════════════════

\begin{frame}{Near-Earth Asteroids and Resupply Channels}
  NEAs ($q < 1.3$~AU, residence time $\sim 10$~Myr) require steady-state resupply from the main belt:
  \vskip4pt
  \begin{itemize}\setlength\itemsep{3pt}
    \item \textbf{Resonances:} Kirkwood (3:1, 5:2) and $\nu_6$ secular resonance pump eccentricity
    \item \textbf{Yarkovsky drift:} thermal recoil drives slow semimajor axis drift, $\dv{a}{t} \propto \dfrac{1}{D \, \rho \, \sqrt{a}}$
    \item \textbf{Hazards:} $\sim$2{,}500 Potentially Hazardous Asteroids ($D > 140$~m, MOID $< 0.05$~AU)
  \end{itemize}
\end{frame}

\begin{frame}
  \ipsherokey[0.65]{neo_sfd_schunova}{Cumulative size-frequency distribution of near-Earth objects against absolute magnitude $H_V$, a power law over five orders of magnitude. Credit: Schunova et al.\ (2017), Fig.~9.}{Impact rates: $10$~m per decade; $50$~m (Tunguska) per $10^3$~yr; $1$~km per $5 \times 10^5$~yr; $10$~km (Chicxulub) per $10^8$~yr.}
\end{frame}

\begin{frame}{Yarkovsky and YORP Detected on Real NEAs}
  \begin{columns}[T]
    \column{0.5\textwidth}
    \includegraphics[width=\linewidth, height=0.52\textheight, keepaspectratio]{yarkovsky_detection_vokrouhlicky}\\
    \footnotesize \emph{Yarkovsky:} (6489) Golevka, Arecibo radar
    \column{0.5\textwidth}
    \includegraphics[width=\linewidth, height=0.52\textheight, keepaspectratio]{yorp_detection_vokrouhlicky}\\
    \footnotesize \emph{YORP:} (54509) YORP rotational acceleration
  \end{columns}
  \vskip6pt
  \keyresult{YORP is the same thermal asymmetry acting on \emph{spin}: it can spin small asteroids to fission or stop them.}
\end{frame}

\begin{frame}
  \ipsherokey[0.61]{dart_lightcurve_daly}{Light curves of Didymos (top, 2.26~h rotation) and Dimorphos (bottom, its new 11.372~h orbit) six days after the DART impact. Thomas et al.\ (2023), Fig.~3.}{The orbital period fell by $33.0 \pm 1.0$ minutes (3$\sigma$) and the momentum enhancement $\beta \approx 3.6$ shows that ejecta recoil exceeds the direct push; Hera (ESA, launched 2024) measures the crater and $\beta$.}
\end{frame}


% ════════════════════════════════════════════════════════════
% BREAK ─ end of inner-solar-system small bodies
% ════════════════════════════════════════════════════════════
% ── Break ───────────────────────────────────────────────────
\breakslide[End of Part 1: Meteorites \& Inner Small Bodies\\Part 2: Trans-Neptunian Region, Comets, Missions]


% ════════════════════════════════════════════════════════════
\sectionimage{ceres_occator}{Occator crater on Ceres, Dawn. NASA/JPL-Caltech}
\section{Dwarf Planets and the Trans-Neptunian Region}
% ════════════════════════════════════════════════════════════

\begin{frame}
  \ipsherokey[0.61]{ceres_occator}{Occator crater on Ceres, Dawn (2015--2018): bright deposits of hydrated $\mathrm{Na_2CO_3}$ and $\mathrm{NH_4Cl}$ from erupted brines. Credit: NASA/JPL-Caltech.}{Ceres ($R = 470$~km, mass $\sim 1\%$ of the Moon) has, or had, a partially liquid water layer at depth, and may be a captured outer-solar-system body (Nice Model: a dynamical instability among the giant planets).}
\end{frame}

\begin{frame}{Trans-Neptunian Region: Five Sub-Populations}
  The orbits preserve the gravitational footprint of giant-planet migration:
  \vskip8pt
  \begin{itemize}\setlength\itemsep{4pt}
    \item \textbf{Classicals ($a = 42$--$47$~AU):} dynamically unperturbed cold classicals ($i \lesssim 5^\circ$) and scattered hot classicals ($i \le 30^\circ$)
    \item \textbf{Resonant (e.g.\ 3:2 plutinos):} bodies swept into mean-motion resonances during Neptune's outward migration
    \item \textbf{Scattered disk and detached:} Neptune-scattered bodies ($q \sim 30$~AU) and decoupled detached objects ($q > 40$~AU)
  \end{itemize}
\end{frame}

\begin{frame}
  \ipsherokey[0.66]{bannister_ossos_orbits}{1142 characterised TNOs from OSSOS in orbital space; the vertical lines mark Neptune's resonances. Bannister et al.\ (2018).}{The dense low-inclination cluster at 42--47 AU is the cold classical belt; the architecture is a fossil of giant-planet migration.}
\end{frame}

\begin{frame}
  \ipsherokey[0.56]{petit_classical_kbo}{The three components of the CFEPS-L7 model of the classical Kuiper belt in $e$ and $i$ against $a$: the hot component, the stirred component, and the dense low-inclination kernel near 44 AU. Petit et al.\ (2011).}{Kernel and hot component share the same $a$ but different $(e, i)$: the cold classicals were never strongly perturbed by Neptune, the hot classicals were excited by the giant-planet instability.}
\end{frame}


% ════════════════════════════════════════════════════════════
\sectionimage{pluto_color_stern}{Pluto, New Horizons (2015). NASA/JHUAPL/SwRI}
\section{Pluto, Charon, Arrokoth}
% ════════════════════════════════════════════════════════════

\begin{frame}
  \ipsherokey[0.61]{pluto_color_stern}{True-colour Pluto from New Horizons Ralph (2015), with Tombaugh Regio and Sputnik Planitia. Credit: NASA/JHUAPL/SwRI.}{With $R = 1187$~km and $\rho = 1860$~kg\,m$^{-3}$: water-ice bedrock under volatile frosts, a convecting $\mathrm{N_2}$-ice basin younger than $10$~Myr, a $\sim 1$~Pa $\mathrm{N_2}$ atmosphere and a probable subsurface ocean.}
\end{frame}

\begin{frame}
  \ipsherokey[0.65]{sputnik_convection_mckinnon}{Numerical model of solid-state convection in Sputnik Planitia at $\mathrm{Ra}_b \sim 3\times 10^5$. McKinnon et al.\ (2016).}{Cells tens of km across overturn in $\sim 10^6$~yr, driven by present-day radiogenic heat in the rocky core.}
\end{frame}

\begin{frame}
  \ipsherokey[0.67]{pluto_charon_features_stern}{Pluto and Charon with their named features, New Horizons (2015). Stern et al.\ (2015).}{Charon ($R = 606$~km, $\sim 1/8$ of Pluto's mass) puts the barycentre \emph{outside} Pluto; the four small moons Styx, Nix, Kerberos and Hydra are likely fragments of an early collisional event.}
\end{frame}

\begin{frame}
  \ipshero{charon_mordor}{Charon from New Horizons: the dark reddish Mordor Macula of tholins from Pluto's escaping atmosphere, and equatorial extensional chasms. Credit: NASA/JHUAPL/SwRI, public domain.}
\end{frame}

\begin{frame}
  \ipsherokey{arrokoth_color_stern}{Arrokoth, a $\sim$35 km \textbf{contact binary}, New Horizons flyby of 1 January 2019. Stern et al.\ (2019).}{Two flat lobes joined at a narrow neck with no impact crater at the join: a merger at $\lesssim$ a few m~s$^{-1}$, formation by streaming-instability gravitational collapse.}
\end{frame}


% ════════════════════════════════════════════════════════════
\sectionimage{sedna_orbits_batygin}{Detached TNO orbits and Planet 9 hypothesis. Batygin et al.\ (2019)}
\section{Other Dwarf Planets, Sedna, and the Oort Cloud}
% ════════════════════════════════════════════════════════════

\begin{frame}{Eris, Haumea, Makemake}
  \begin{columns}[T]
    \column{0.33\textwidth}
    \includegraphics[width=\linewidth, height=0.45\textheight, keepaspectratio]{eris}\\
    \footnotesize \emph{Eris} + Dysnomia. More massive than Pluto.
    \column{0.34\textwidth}
    \includegraphics[width=\linewidth, height=0.45\textheight, keepaspectratio]{haumea_ring}\\
    \footnotesize \emph{Haumea} occultation: revealed a \textbf{ring}.
    \column{0.33\textwidth}
    \includegraphics[width=\linewidth, height=0.45\textheight, keepaspectratio]{makemake}\\
    \footnotesize \emph{Makemake} + small moon. Methane-ice surface.
  \end{columns}
  \vskip4pt\relax
  Plus dwarf-planet candidates: Gonggong, Quaoar, Orcus, Salacia, Sedna.
\end{frame}

\begin{frame}
  \ipsherokey[0.66]{sedna_orbits_batygin}{Orbits of the detached TNOs and the Planet 9 hypothesis. Batygin et al.\ (2019).}{Sedna ($q = 76$~AU, $Q = 900$~AU, period 11{,}400~yr) is too far for Neptune and too close for present-day stars: an inner Oort cloud or birth-cluster relic; the orbital clustering motivates the \textbf{Planet 9} hypothesis.}
\end{frame}

\begin{frame}
  \ipsherokey[0.6]{oort_cloud_kaib}{Two simulations of Oort cloud formation: perihelion $q$ (top) and inclination (bottom) against semimajor axis $a$. Kaib \& Quinn (2008).}{$2{,}000$--$50{,}000$ AU, roughly spherical, $\sim 10^{11}$--$10^{12}$ objects of $\sim 1$--$10\,\Mearth$ in total, planetesimals scattered from the Jupiter-Neptune zone and lifted by the galactic tide and passing stars; inferred by Oort in 1950, \textbf{never directly observed}.}
\end{frame}


% ════════════════════════════════════════════════════════════
\sectionimage{halley}{Comet 1P/Halley nucleus, Giotto (1986). ESA/MPAe Lindau}
\section{Comets: Structure, D/H, and Origin}
% ════════════════════════════════════════════════════════════

\begin{frame}
  \ipsherokey[0.6]{comet_diagram}{Anatomy of an active comet: nucleus, coma, ion tail and dust tail. Credit: Wikimedia Commons, public domain.}{Whipple's dirty snowball: a dark, porous \textbf{nucleus} of $1$--$30$~km (67P: $533$ kg m$^{-3}$, $\sim 70\%$ porosity) of ices, dust and organics; inside $\sim 5$~AU water sublimation feeds a $10^4$--$10^6$~km \textbf{coma}, a straight anti-sunward \textbf{ion tail} and a curved \textbf{dust tail}.}
\end{frame}

\begin{frame}
  \ipsherokey[0.65]{halley}{Halley's nucleus ($15 \times 8 \times 8$~km) with jets from small active patches, Giotto Halley Multicolour Camera, 13--14 March 1986. Credit: ESA/MPAe Lindau.}{A $76$-year retrograde orbit with $T_J \approx -0.6$: an Oort-cloud comet captured to a short period.}
\end{frame}

\begin{frame}{Short-Period vs Long-Period: The Tisserand Parameter}
  \[
    T_J = \frac{a_J}{a} + 2 \sqrt{\frac{a}{a_J}\,(1 - e^2)} \cos i, \qquad a_J = 5.20\,\mathrm{AU}
  \]
  \vskip4pt
  A conserved dynamical invariant that separates asteroids from Jupiter-family and long-period comets:
  \vskip4pt
  \begin{itemize}\setlength\itemsep{3pt}
    \item \textbf{Asteroids ($T_J > 3$):} main-belt orbits decoupled from Jupiter encounters (e.g.\ Ceres, $T_J \approx 3.3$)
    \item \textbf{Jupiter-family comets ($2 < T_J < 3$):} Kuiper-belt-derived comets controlled by Jupiter (e.g.\ 67P, $T_J \approx 2.7$)
    \item \textbf{Long-period comets ($T_J < 2$):} Oort-cloud visitors on high-inclination orbits; Halley-type comets ($T_J \approx -0.6$) are Oort-derived but captured to short orbits (76 yr)
  \end{itemize}
\end{frame}

\begin{frame}
  \ipshero[0.76]{67p_nucleus}{Bilobed nucleus of Jupiter-family comet 67P/Churyumov-Gerasimenko ($T_J \approx 2.7$), two lobes joined at a narrow neck by a gentle primordial merger, Rosetta OSIRIS. Credit: ESA/Rosetta/MPS for OSIRIS Team.}
\end{frame}

\begin{frame}
  \ipsherokey[0.6]{hartley2}{Comet 103P/Hartley 2 from EPOXI (2010), whose water D/H matches Earth's oceans. Credit: NASA/JPL-Caltech/UMD, public domain.}{Earth's oceans have D/H $= 1.56 \times 10^{-4}$ (VSMOW); Hartley 2 matches, 67P is $\approx 3\times$ and long-period comets $\approx 2\times$ VSMOW, so no comet class matches Earth: carbonaceous chondrites dominate Earth's water and comets supply at most a minor fraction.}
\end{frame}


% ════════════════════════════════════════════════════════════
\sectionimage{67p_nucleus}{Comet 67P/Churyumov-Gerasimenko, Rosetta. ESA/Rosetta/MPS for OSIRIS Team}
\section{Recent Missions and Sample Return}
% ════════════════════════════════════════════════════════════

\begin{frame}{Three Sample Returns: Itokawa, Ryugu, Bennu}
  \begin{columns}[T]
    \column{0.33\textwidth}
    \includegraphics[width=\linewidth, height=0.4\textheight, keepaspectratio]{itokawa_full}\\
    \footnotesize \emph{Hayabusa, 2010}: 1{,}500 grains. Itokawa $=$ LL chondrite parent.
    \column{0.34\textwidth}
    \includegraphics[width=\linewidth, height=0.4\textheight, keepaspectratio]{ryugu}\\
    \footnotesize \emph{Hayabusa2, 2020}: 5.4 g. Ryugu $\approx$ CI chondrite. Amino acids, nucleobases.
    \column{0.33\textwidth}
    \includegraphics[width=\linewidth, height=0.4\textheight, keepaspectratio]{bennu_lauretta}\\
    \footnotesize \emph{OSIRIS-REx, 2023}: 121 g. Bennu CM/CI-like, Mg carbonate veins.
  \end{columns}
  \vskip4pt
  \keyresult{Each rock $\to$ specific parent body with a known orbit and full in-situ context.}
\end{frame}

\begin{frame}
  \ipshero[0.76]{dinkinesh_levison}{The Dinkinesh-Selam system from the Lucy flyby of 1 November 2023: the $\sim$720 m main-belt asteroid with its equatorial ridge, and the moonlet Selam, a contact binary of two $\sim$210 and 230 m lobes. Reproduced from Levison et al.\ (2024).}
\end{frame}

\begin{frame}
  \ipshero{psyche_shape_shepard}{The shape of (16) Psyche from radar and adaptive-optics data, seen from above its south pole, with the regions Alpha, Bravo and Charlie where the body falls short of the best-fit ellipsoid. Reproduced from Shepard et al.\ (2021).}
\end{frame}


% ════════════════════════════════════════════════════════════
\sectionimage{oumuamua_discovery_meech}{1I/'Oumuamua deep image. ESO / K.\ Meech et al.}
\section{Interstellar Visitors}
% ════════════════════════════════════════════════════════════

\begin{frame}
  \ipsherokey[0.6]{oumuamua_discovery_meech}{Deep image of 1I/'Oumuamua at the centre; the stars trail because the telescope tracked the target. Credit: ESO / K.\ Meech et al.}{No detectable coma at the limit of stacked imaging, yet the outgoing trajectory shows a non-gravitational acceleration; its nature is still debated: H$_2$ ice fragment, N$_2$ ice, or tidal disruption shard.}
\end{frame}

\begin{frame}
  \ipsherokey[0.61]{oumuamua_iso_density}{Implied number density of interstellar objects from the single Pan-STARRS detection of 'Oumuamua. 'Oumuamua ISSI Team (2019).}{One detection implies $\sim 10^{4}$ ISOs larger than 100 m within Neptune's orbit \emph{at any moment}, each a sample of a different planetary system; every star ejects planetesimals at rates similar to ours.}
\end{frame}

\begin{frame}
  \ipshero[0.76]{borisov_composition_bodewits}{Volatile composition in the coma of interstellar comet 2I/Borisov against solar-system comets: elevated C/O and depleted H/O point to formation beyond the CO snowline of its host system. Credit: Bodewits et al.\ (2020), Fig.~3.}
\end{frame}


% ════════════════════════════════════════════════════════════
\section*{Summary}
% ════════════════════════════════════════════════════════════

\begin{frame}{Summary}
  \begin{enumerate}\setlength\itemsep{6pt}
    \item Small bodies are \textbf{formation fossils}; \textbf{Pb-Pb dating} anchors the solar system's absolute clock at $4567.30 \pm 0.16$~Myr
    \item The \textbf{NC-CC dichotomy} separates inner and outer reservoirs for $\sim 3$--$4$~Myr; resonance gaps and Yarkovsky thermal recoil drive steady-state NEA delivery
    \item Cometary D/H shows carbonaceous chondrites delivered most of Earth's water
  \end{enumerate}
  \vskip6pt
  \textbf{Next: Lecture 13} turns to the exoplanet revolution.
\end{frame}

% ── Final slide ─────────────────────────────────────────────
\begin{frame}[plain,noframenumbering]
  \begin{tikzpicture}[remember picture, overlay]
    \ipscontain{title_background}
    \fill[ipsPrimary, opacity=0.70] (current page.north west) rectangle (current page.south east);
  \end{tikzpicture}
  \vfill
  \begin{center}
    {\color{white}\Huge \textbf{Questions?}}\\[14pt]
    {\color{white}\large Next: \textbf{Lecture 13. Exoplanets}}\\[10pt]
    {\color{white}\normalsize Lecture notes: \texttt{formingworlds.github.io/IntroductionToPlanetaryScience}}
  \end{center}
  \vfill
\end{frame}

\end{document}
